\subsection{决策单调性与二分队列}
考虑一类一维分段 DP
\[
    f_i=\min_{0\le j<i}\{f_j+w(j+1,i)\}.
\]
若转移矩阵是 Monge 的，即对 $j<J\le i<I$ 有四边形不等式
\[
    w(j+1,i)+w(J+1,I)\le w(j+1,I)+w(J+1,i),
\]
则最优决策点 $p_i=\operatorname*{argmin}_j\{f_j+w(j+1,i)\}$ 单调不减。加上只与 $j$ 有关的 $f_j$ 不会破坏 Monge 性。

对区间费用 $C(l,r)$，常写成
\[
    C(a,c)+C(b,d)\le C(a,d)+C(b,c),\qquad a\le b\le c\le d.
\]
在完整的二维定义域上，也可以只检查相邻的 $2\times2$ 条件
\[
    C(l,r)+C(l+1,r+1)\le C(l,r+1)+C(l+1,r),
\]
它等价于所有四边形不等式。实际题目中通常直接从费用的结构证明；例如有序点分段后到中位数的 $L_1$ 代价就是常见的 Monge 费用。

二分队列维护若干候选决策 $(j,L)$，表示从位置 $L$ 开始，$j$ 是当前队列中最优的决策。由于决策单调，两个候选的优劣只会发生一次切换：插入新决策 $i$ 时，与队尾旧决策比较，并二分出新决策第一次不劣的位置；若新决策从队尾的起点起就更优，则直接弹掉队尾。查询 $f_i$ 时只需根据这些分界点移动队首。每个决策至多入队、出队一次，另有一次二分，因此总复杂度 $O(n\log n)$。

下面是最小化问题的核心模板。\texttt{better(j,J,x)} 表示在终点 $x$ 处决策 $j$ 是否严格优于 $J$；相等时的 tie-break 要与 WQS 中记录的段数保持一致。
\begin{lstlisting}
ll dp[N];
int cnt[N];
struct Node{ int j,l; } q[N];
int hd,tl;

bool better(int j,int J,int i){
    ll x=dp[j]+cost(j+1,i);
    ll y=dp[J]+cost(J+1,i);
    if(x!=y)return x<y;
    return cnt[j]<cnt[J];
}

void take(int j,int i,ll extra=0){
    dp[i]=dp[j]+cost(j+1,i)+extra;
    cnt[i]=cnt[j]+1;
}

void insert_candidate(int i){
    if(i>=n)return;
    while(hd<=tl){
        int j=q[tl].j;
        int l=max(i+1,q[tl].l),r=n;

        // old candidate is still better at n: i never becomes optimal
        if(better(j,i,n))return;

        // first position where i is no worse than j
        while(l<r){
            int mid=(l+r)>>1;
            if(better(j,i,mid))l=mid+1;
            else r=mid;
        }

        if(l==q[tl].l)--tl;          // i dominates the whole last interval
        else{
            q[++tl]={i,l};
            return;
        }
    }
    q[++tl]={i,i+1};
}

void solve_dp(ll extra=0){
    dp[0]=0;cnt[0]=0;
    hd=tl=1;q[1]={0,1};
    for(int i=1;i<=n;i++){
        while(hd<tl&&q[hd+1].l<=i)++hd;
        take(q[hd].j,i,extra);
        insert_candidate(i);
    }
}
\end{lstlisting}

和 WQS 合用时，令 \texttt{extra=k}，一次 \texttt{check(k)} 就是一次 $O(n\log n)$ 的上述 DP，于是总复杂度再乘一个 WQS 的二分次数。


\paragraph{使用条件与多测} solve\_dp 每次重置 dp[0]、cnt[0]、队列边界；其余 dp/cnt 在使用前被顺序覆盖。cost 的预处理由调用方重建。Monge 条件须证明且两候选含 tie-break 的比较须单次切换；和 WQS 组合时也应验证段数 tie-break 满足该条件，不能只验证数值最优点单调。要求 cost 为 $O(1)$ 才有所述复杂度；中间加法须防溢出。
